\SourceLeaf{235}%
adeoque
\[
\begin{aligned}
\frac{\mathrm{d}X}{\mathrm{d}x}&=Z^3(-(1+uu)T+tuU)\\
\frac{\mathrm{d}X}{\mathrm{d}y}&=Z^3(-(1+uu)U+tuV)\\
\frac{\mathrm{d}Y}{\mathrm{d}x}&=Z^3(tuT-(1+tt)U)\\
\frac{\mathrm{d}Y}{\mathrm{d}y}&=Z^3(tuU-(1+tt)V)
\end{aligned}
\]
quibus valoribus in expressione praecedente substitutis, prodit
\[
k=Z^6(TV-UU)(1+tt+uu)=Z^4(TV-UU)=\frac{TV-UU}{(1+tt+uu)^2}
\]
\Article{8}
Per idoneam electionem initii et axium coordinatarum facile effici potest, ut pro puncto determinato \(A\) valores quantitatum \(t,u,U\) evanescant. Scilicet duae priores conditiones iam adimplentur, si planum tangens in hoc puncto pro plano coordinatarum \(x,y\) adoptatur. Quarum initium si insuper in puncto \(A\) ipso collocatur, manifesto expressio coordinatarum \(z\) adipiscitur formam talem
\[
z=\tfrac{1}{2}T^0xx+U^0xy+\tfrac{1}{2}V^0yy+\Omega
\]
ubi \(\Omega\) erit ordinis altioris quam secundi. Mutando dein situm axium ipsarum \(x,y\) angulo \(M\) tali ut habeatur
\[
\mathrm{tang}\,2M=\frac{2U^0}{T^0-V^0}
\]
facile perspicitur, prodituram esse aequationem huius formae
\[
z=\tfrac{1}{2}Txx+\tfrac{1}{2}Vyy+\Omega
\]
quo pacto etiam tertiae conditioni satisfactum est. Quibus ita factis, patet

I. Si superficies curva secetur plano ipsi normali et per axem coordinatarum \(x\) transeunte, oriri curvam planam, cuius radius curvaturae in puncto \(A\) fiat \(=\frac{1}{T}\), signo positivo vel negativo indicante concavitatem vel convexitatem versus plagam eam, versus quam coordinatae \(z\) sunt positivae.

II. Simili modo \(\frac{1}{V}\) erit in puncto \(A\) radius curvaturae curvae planae, quae oritur per sectionem superficiei curvae cum plano per axes ipsarum \(y,z\) transeunte.
